\documentclass[12pt]{article}

% ==============================================================================
% PREÂMBULO: PACOTES E CONFIGURAÇÕES
% ==============================================================================
% Utilizamos o template da SBC (Sociedade Brasileira de Computação) para manter
% a formatação acadêmica padrão de eventos como o CSBC.
\usepackage{sbc-template}

% Gerenciamento de imagens e URLs.
\usepackage{graphicx,url}

% Codificação UTF-8 para caracteres acentuados e idioma português do Brasil.
\usepackage[utf8]{inputenc}
\usepackage[brazil]{babel}

% Melhora a aparência das tabelas com linhas horizontais de alta qualidade.
\usepackage{booktabs}

% Ambiente matemático para fórmulas (ex.: complexidade O(N×M)).
\usepackage{amsmath}
\usepackage{amsfonts}

% Corrige um conflito específico do BibTeX no Linux Mint (mantido conforme original).
\let\em\itshape 

% Permite que o LaTeX quebre linhas de maneira mais flexível, evitando overfull boxes.
\sloppy

% ==============================================================================
% INFORMAÇÕES DO ARTIGO
% ==============================================================================
% O título reflete a comparação central do trabalho: precisão exata vs. eficiência.
\title{Análise Comparativa de Algoritmos para Alinhamento de Sequências de DNA: Programação Dinâmica vs. Abordagem Gulosa}

% Autores do trabalho (preencha com os nomes completos da equipe).
\author{
    Rafael Passos Domingues\inst{1}, 
    Wagner Donizete Gonçalves\inst{1}
}

% Afiliação institucional e e-mail de contato do primeiro autor.
\address{
    Departamento de Ciência da Computação -- Universidade Federal de Alfenas (UNIFAL-MG)\\
    Avenida Jovino Fernandes de Sales 2600 -- CEP 37133-840 -- Alfenas -- MG -- Brasil
    \email{rafael.domingues@sou.unifal-mg.edu.br, wagner.goncalves@sou.unifal-mg.edu.br}
}

\begin{document}

\maketitle

% ==============================================================================
% RESUMO
% ==============================================================================
% O resumo condensa: problema, métodos, métricas (score biológico), e a conclusão
% principal sobre o trade-off e a limitação de hardware.
\begin{resumo}
    Este trabalho apresenta uma análise comparativa entre duas abordagens algorítmicas aplicadas ao problema de alinhamento global de sequências de DNA: a Programação Dinâmica (baseada no algoritmo de Needleman-Wunsch) e uma heurística Gulosa. O sistema de pontuação adotado modela fielmente a complementaridade biológica das bases nitrogenadas, onde um \textit{match} verdadeiro exige os pares Adenina-Timina ou Citosina-Guanina. Os resultados experimentais evidenciam o clássico \textit{trade-off} entre precisão e eficiência computacional. Enquanto a Programação Dinâmica garante a otimalidade global com custo quadrático $O(N \times M)$, o algoritmo Guloso opera em tempo linear $O(N)$, porém com perda significativa de qualidade no alinhamento. A análise de desempenho revela que, para sequências superiores a 10.000 bases, a solução exata torna-se proibitiva devido a limitações de memória RAM e \textit{cache misses}, justificando a adoção de heurísticas em aplicações de bioinformática em larga escala.
\end{resumo}

% ==============================================================================
% 1. INTRODUÇÃO
% ==============================================================================
% A introdução estabelece a ponte entre biologia e computação, explicando por que
% o alinhamento de DNA não é uma simples comparação de strings.
\section{Introdução}
\label{sec:introducao}

Desde a elucidação da estrutura em dupla hélice do DNA por \cite{watson:53}, o alinhamento de sequências genéticas tornou-se uma ferramenta fundamental na bioinformática. O DNA pode ser interpretado como a linguagem de programação mais antiga do planeta: um código quaternário composto pelas bases nitrogenadas Adenina (A), Timina (T), Citosina (C) e Guanina (G). Diferentemente de uma comparação textual convencional, onde letras iguais indicam correspondência, o pareamento do DNA obedece a regras bioquímicas estritas de complementaridade: a Adenina liga-se exclusivamente à Timina (A-T), e a Citosina à Guanina (C-G), formando pontes de hidrogênio estáveis.

O problema do alinhamento global consiste em encontrar a melhor sobreposição entre duas cadeias, inserindo espaços (\textit{gaps}) quando necessário, de modo a maximizar uma função de pontuação que reflita a probabilidade evolutiva daquele alinhamento. Computacionalmente, isso se traduz em um problema de otimização combinatória com vasto espaço de busca.

Neste projeto, implementamos e comparamos dois paradigmas algorítmicos distintos:
\begin{itemize}
    \item \textbf{Programação Dinâmica (PD)}: Baseada no algoritmo de \cite{needleman:70}, garante a solução ótima global ao custo de maior consumo de tempo e memória.
    \item \textbf{Algoritmo Guloso (Heurística)}: Toma decisões locais e imediatistas, priorizando a velocidade de execução em detrimento da garantia de otimalidade.
\end{itemize}

O foco deste trabalho não se limita à corretude matemática dos algoritmos, mas estende-se à análise de seu comportamento sob as restrições físicas do hardware, como latência de memória, fragmentação e \textit{cache misses}, tópicos essenciais na engenharia de software de alto desempenho \cite{cormen:2012,ziviani:2007}.

% ==============================================================================
% 2. METODOLOGIA
% ==============================================================================
% Detalha a função de avaliação biológica e a implementação de ambos os algoritmos,
% incluindo as otimizações de baixo nível na alocação de memória.
\section{Metodologia}
\label{sec:metodologia}

\subsection{Sistema de Pontuação Biológico}
\label{subsec:pontuacao}

A modelagem matemática do problema exige uma função de \textit{score} que traduza a química molecular em valores numéricos. Diferente de algoritmos de comparação de \textit{strings} que recompensam caracteres idênticos, nosso sistema segue rigorosamente os critérios de complementaridade de Watson e Crick. A Tabela~\ref{tab:tabela_pontuacao} define os valores utilizados.

\begin{table}[!htb]
\centering
\caption{Função de Avaliação Baseada em Complementaridade Biológica}
\label{tab:tabela_pontuacao}
\begin{tabular}{@{}llc@{}}
\toprule
\textbf{Tipo de Alinhamento} & \textbf{Exemplo} & \textbf{Pontuação} \\ \midrule
\textit{Match} (Par Complementar)       & A-T, C-G         & +2 \\
\textit{Mismatch} (Não Complementar)    & A-A, A-C, T-G    & -1 \\
\textit{Gap} (Inserção/Deleção)         & A-, -C           & -2 \\ \bottomrule
\end{tabular}
\end{table}

\textbf{Justificativa Biológica:}
\begin{itemize}
    \item \textbf{Match (+2)}: Representa a conservação evolutiva. A recompensa é alta porque o encaixe químico perfeito preserva a função da proteína codificada.
    \item \textbf{Mismatch (-1)}: Corresponde a uma mutação pontual (substituição). Embora possa alterar a estrutura proteica, a célula frequentemente tolera esses erros. A penalidade é moderada.
    \item \textbf{Gap (-2)}: Simula uma mutação do tipo \textit{indel} (inserção ou deleção). Quando uma base é removida ou adicionada, toda a matriz de leitura do código genético (que é lido em tríades) se desloca (\textit{frameshift}), geralmente resultando em uma proteína não funcional. Por isso, a penalidade é severa.
\end{itemize}

Note que, em nossa implementação, a igualdade entre caracteres (ex.: \texttt{'A'} com \texttt{'A'}) é tratada como \textit{mismatch} e penalizada com -1, pois a natureza não recompensa a igualdade, mas sim a complementaridade.

\subsection{Programação Dinâmica (Needleman-Wunsch)}
\label{subsec:pd}

A Programação Dinâmica resolve o problema dividindo-o em subproblemas menores e armazenando os resultados intermediários. Construímos uma matriz de dimensão $(N+1) \times (M+1)$, onde $N$ e $M$ são os comprimentos das sequências. Cada célula $(i,j)$ armazena a pontuação máxima do alinhamento parcial entre os primeiros $i$ caracteres da sequência 1 e os primeiros $j$ caracteres da sequência 2.

O preenchimento segue a recorrência:
\[
M[i][j] = \max
\begin{cases}
M[i-1][j-1] + \text{score}(S_1[i], S_2[j]) & \text{(alinhar diagonal)} \\
M[i-1][j] + \text{gap\_penalty} & \text{(gap na sequência 2)} \\
M[i][j-1] + \text{gap\_penalty} & \text{(gap na sequência 1)}
\end{cases}
\]

\textbf{Otimização de Engenharia de Software (Baixo Nível):}
% ------------------------------------------------------------------------------
% COMENTÁRIO EM BLOCO: EXPLICAÇÃO DA OTIMIZAÇÃO DE MEMÓRIA
% ------------------------------------------------------------------------------
% Em vez de alocar uma matriz bidimensional clássica (array de ponteiros para
% arrays), que causa fragmentação e perda de localidade espacial, alocamos um
% ÚNICO bloco contíguo de memória de tamanho (N+1)*(M+1). O acesso é feito via
% aritmética de ponteiros: índice = i * cols + j.
% Essa técnica reduz drasticamente os "cache misses" porque a CPU consegue
% pré-carregar linhas inteiras da matriz na cache L1/L2. Em alinhamentos de
% sequências longas, essa diferença pode representar uma economia de 30% a 50%
% no tempo de execução.
% ------------------------------------------------------------------------------
Para otimizar o uso da hierarquia de memória e evitar \textit{cache misses} frequentes, não utilizamos um array bidimensional tradicional. Em vez disso, alocamos um vetor unidimensional contíguo utilizando \texttt{malloc((N+1) * (M+1) * sizeof(int))}. O acesso à posição $(i,j)$ é realizado através da aritmética de ponteiros \texttt{vetor[i * cols + j]}. Essa abordagem garante localidade espacial, pois os dados de uma mesma linha residem em endereços físicos adjacentes, maximizando a eficiência da cache do processador \cite{kernighan:89}.

A complexidade de tempo e espaço deste algoritmo é $O(N \times M)$. Embora garanta a solução ótima, sua viabilidade é limitada pelo hardware para sequências longas.

\subsection{Algoritmo Guloso}
\label{subsec:guloso}

Em contraste com a abordagem exaustiva da Programação Dinâmica, o algoritmo Guloso implementa uma heurística de decisão local. A cada passo, ele examina apenas os caracteres atuais e, no máximo, uma posição à frente (\textit{lookahead} de profundidade 1).

A lógica de decisão é a seguinte:
\begin{enumerate}
    \item Se os caracteres atuais são complementares (A-T ou C-G), alinha-os diretamente (ganho +2) e avança ambos os índices.
    \item Caso contrário, o algoritmo verifica se inserir um \textit{gap} em uma das sequências permitirá um \textit{match} na posição seguinte.
    \item Se nenhuma dessas opções for promissora, aceita-se o \textit{mismatch} atual (penalidade -1).
    \item Ao final do processamento, eventuais caracteres restantes são tratados como \textit{gaps} forçados.
\end{enumerate}

Essa estratégia evita a construção de matrizes e percorre as sequências de forma linear, resultando em complexidade de tempo $O(N + M)$. O consumo de memória é constante, utilizando apenas algumas variáveis escalares e ponteiros.

Contudo, por se basear em escolhas míopes (ótimo local), o algoritmo não possui garantia de convergência para a solução global ótima. Em muitos casos, o acúmulo de pequenas decisões subótimas leva a alinhamentos com pontuação significativamente inferior.

% ==============================================================================
% 3. RESULTADOS E DISCUSSÃO
% ==============================================================================
% Apresenta os dados experimentais (tabelas e gráficos) e interpreta o trade-off
% à luz das limitações de hardware e das dimensões do genoma real.
\section{Resultados e Discussão}
\label{sec:resultados}

Os experimentos foram conduzidos em um ambiente controlado, com múltiplas execuções para cada tamanho de entrada, a fim de mitigar flutuações de desempenho do sistema operacional. A Tabela~\ref{tab:comparacao_scores} apresenta a média das pontuações obtidas pelos dois algoritmos.

\begin{table}[!htb]
\centering
\caption{Comparação da Qualidade do Alinhamento (Score Médio)}
\label{tab:comparacao_scores}
\begin{tabular}{@{}lccc@{}}
\toprule
\textbf{Tamanho da Sequência (N)} & \textbf{Score DP (Ótimo)} & \textbf{Score Guloso} & \textbf{Diferença Absoluta} \\ \midrule
100   & 35    & -1    & 36    \\
500   & 196   & -44   & 240   \\
1.000 & 408   & -43   & 451   \\
5.000 & 2.174 & -179  & 2.353 \\
10.000& 4.443 & -450  & 4.893 \\ \bottomrule
\end{tabular}
\end{table}

\textbf{Análise da Qualidade:}
Observa-se que a Programação Dinâmica mantém pontuações consistentemente positivas e crescentes, refletindo sua capacidade de encontrar o alinhamento evolutivamente mais plausível. Por outro lado, o algoritmo Guloso frequentemente atinge escores negativos, indicando que as penalidades acumuladas superam os acertos. Isso ocorre porque as decisões locais do Guloso, embora rápidas, carecem de uma visão global do contexto da sequência.

A Tabela~\ref{tab:resultados_tempo} quantifica o custo computacional dessa precisão.

\begin{table}[!htb]
\centering
\caption{Benchmarking de Desempenho (Tempo Médio em Segundos)}
\label{tab:resultados_tempo}
\begin{tabular}{@{}lcc@{}}
\toprule
\textbf{Tamanho (N)} & \textbf{Tempo DP (s)} & \textbf{Tempo Guloso (s)} \\ \midrule
1.000  & 0.025 & 0.001 \\
5.000  & 0.501 & 0.005 \\
10.000 & 2.004 & 0.010 \\ \bottomrule
\end{tabular}
\end{table}

A Programação Dinâmica exibe um crescimento quadrático no tempo de execução, conforme esperado pela complexidade $O(N \times M)$. O algoritmo Guloso, por sua vez, apresenta um crescimento linear quase imperceptível na escala do gráfico.

\subsection{Análise Gráfica e o Limite do Hardware}
\label{subsec:graficos}

A Figura~\ref{fig:grafico_desempenho} consolida visualmente o \textit{trade-off} entre desempenho e consumo de memória.

\begin{figure}[!htb]
\centering
% Substitua 'desempenho_algoritmos_pb' pelo nome real do seu arquivo de imagem (ex: grafico.png)
\includegraphics[width=0.85\textwidth]{desempenho_algoritmos_pb}
\caption{Comparativo de Desempenho: Tempo de Execução e Consumo de Memória}
\label{fig:grafico_desempenho}
\end{figure}

O gráfico de tempo (esquerda) evidencia a trajetória parabólica da Programação Dinâmica, que rapidamente se distancia do eixo horizontal. O gráfico de memória (direita, em escala logarítmica) revela o verdadeiro gargalo: a alocação da matriz de Programação Dinâmica consome toda a memória RAM disponível para sequências acima de 10.000 bases.

\textbf{O Contexto do Genoma Real:}
Para dimensionar a gravidade desse gargalo, considere as dimensões reais do código genético:
\begin{itemize}
    \item \textbf{Vírus SARS-CoV-2}: $\approx 30.000$ bases.
    \item \textbf{Bactéria \textit{E. coli}}: $\approx 4,6 \times 10^6$ bases.
    \item \textbf{Ser Humano (\textit{Homo sapiens})}: $\approx 3,2 \times 10^9$ bases.
\end{itemize}

Se a Programação Dinâmica já se mostra inviável para alinhar duas sequências de 10.000 caracteres (exigindo $\approx 100$ milhões de células na matriz), a tentativa de alinhar dois genomas humanos demandaria uma matriz de aproximadamente $10^{19}$ células. Isso está ordens de grandeza além da capacidade de memória de qualquer supercomputador existente. É fisicamente impossível executar o algoritmo exato de Needleman-Wunsch em escala genômica.

Essa constatação valida a existência e a importância de algoritmos heurísticos como o Guloso. Embora ele não encontre o alinhamento perfeito definido por Watson e Crick em todos os detalhes, ele fornece uma resposta ``boa o suficiente'' em tempo hábil, permitindo que a pesquisa biológica avance dentro das limitações do hardware atual.

% ==============================================================================
% 4. DIVISÃO DE TAREFAS
% ==============================================================================
% Seção opcional para relatórios de projeto em grupo.
\section{Divisão de Tarefas}
\begin{itemize}
    \item \textbf{Equipe Alpha (Rafael e Wagner):} Projeto arquitetural do código em C, implementação das otimizações de memória contígua e análise de \textit{cache misses}.
    \item \textbf{Equipe Bravo (Terceiro e Quarto):} Geração de instâncias de teste, coleta de dados em CSV e execução dos benchmarks de tempo.
    \item \textbf{Equipe Charlie (Quinto e Sexto):} Análise estatística dos resultados, elaboração dos gráficos e redação do presente relatório.
\end{itemize}

% ==============================================================================
% 5. CONCLUSÕES
% ==============================================================================
\section{Conclusões}
\label{sec:conclusoes}

Este trabalho demonstrou empiricamente o \textit{trade-off} inerente ao projeto de algoritmos para alinhamento de sequências biológicas. A Programação Dinâmica oferece a solução ótima global, mas seu custo computacional $O(N \times M)$ a torna inaplicável para problemas de escala genômica devido a limitações físicas de memória e tempo. Em contrapartida, o Algoritmo Guloso, com complexidade linear $O(N)$, sacrifica a garantia de otimalidade para viabilizar a execução em grandes volumes de dados.

A principal contribuição deste estudo para a formação em Ciência da Computação é a percepção de que a teoria algorítmica não pode ser dissociada das restrições da arquitetura de computadores. A otimização de baixo nível (como a alocação contígua de memória) e a compreensão do comportamento da cache são tão cruciais quanto a corretude matemática do algoritmo.

Em um cenário futuro onde a computação quântica ou novas arquiteturas de memória superem essas barreiras, a lógica fundamental do alinhamento de sequências -- os \textit{matches}, \textit{mismatches} e \textit{gaps} definidos por Watson e Crick -- permanecerá a mesma. Até lá, a bioinformática continuará a depender do equilíbrio inteligente entre a precisão da Programação Dinâmica e a velocidade das heurísticas.

% ==============================================================================
% REFERÊNCIAS BIBLIOGRÁFICAS
% ==============================================================================
% O estilo 'sbc' formata as referências conforme as normas da SBC.
\bibliographystyle{sbc}
% O arquivo .bib deve estar no mesmo diretório e conter as entradas citadas.
\bibliography{sbc-template}

\end{document}